\documentclass[11pt]{article}
\pdfinfoomitdate=1\pdftrailerid{}
\usepackage[letterpaper,margin=.7in,headheight=15pt,headsep=18pt,footskip=24pt]{geometry}
\usepackage[T1]{fontenc}\usepackage{lmodern,amsmath,amssymb,array,booktabs,fancyhdr,needspace,hyperref,tikz}
\usetikzlibrary{arrows.meta}
\renewcommand{\familydefault}{\sfdefault}\setlength{\parindent}{0pt}\setlength{\parskip}{5pt}
\pagestyle{fancy}\fancyhf{}\fancyhead[L]{\small Week NN / Untangling networks / Facilitator guide}
\fancyfoot[L]{\scriptsize Bellingham Math Circle / Week NN / FNN-FAC-v1}\fancyfoot[R]{\thepage}
\newcommand{\head}[1]{\par\needspace{5\baselineskip}\textbf{\large #1}\par\nopagebreak}
\newcommand{\sub}[1]{\par\needspace{3\baselineskip}\textbf{#1}\par\nopagebreak}
\newcommand{\fig}[2]{\begin{tabular}[b]{@{}c@{}}\input{gen/#1}\\[2pt]{\scriptsize #2}\end{tabular}}
\begin{document}
\head{Mathematical overview: counting regions proves a drawing impossible}
\textbf{The model.} A network is dots joined by links, never two links between the same two dots. Dots move and links bend, but a link never lies across a dot. A crossing is a place where exactly two links cross; the 2--3 and 4--5 rules forbid more than two links at one place, which the six-dot lower bound needs (see Fewest crossings). K--1 needs no such rule. A network is \emph{untangled} when it is drawn with no crossings (it is \emph{planar}). By F\'ary's theorem, anything that untangles with bent links also untangles with straight ones.

\textbf{Euler's rule.} An untangled network in one piece with $V$ dots, $E$ links and $F$ regions (the outside counted) has $V-E+F=2$. \emph{Proof idea:} grow the drawing link by link from one dot ($1-0+1=2$), keeping it in one piece. A link to a new dot adds a dot and a link. A link between two dots already there closes a loop, which cuts one region in two, adding a link and a region. Neither step changes $V-E+F$. \emph{Limits:} $k$ separate pieces give $1+k$ (two triangles: $6-6+3=3$). Every untangled drawing of one network has the same number of regions. A drawing in one piece with $c$ crossings has $E-V+2+c$ regions, since each crossing acts as an extra dot.

\textbf{At most $3V-6$ links.} In an untangled network in one piece with at least 3 dots, every region has at least 3 link-sides round it, and each link has two sides, so $3F\le 2E$; with Euler's rule, $E\le 3V-6$. With no triangle (no three dots all joined to each other) each region needs at least 4 link-sides, and $E\le 2V-4$. A network in several pieces can be joined up first, which only adds links.

\textbf{Two networks can never be untangled.} Five dots all joined ($K_5$: network 2, K--1 Problem 5) has $10>3\cdot5-6$ links. Three houses each joined to three wells ($K_{3,3}$: network 5, K--1 Problem 3) has no triangle, since every link joins a house to a well, and $9>2\cdot6-4$. Two proofs without the formula use only that a closed loop has an inside and an outside. In $K_5$ the ring A--B--C--D--E is a loop; the five other links each run inside or outside it, and only two fit on each side without crossing. In $K_{3,3}$ the ring A--D--B--E--C--F is a loop; A--E, B--F and C--D join opposite corners, so any two on the same side cross, and three links cannot use two sides. A third proof for $K_5$: four dots all joined untangle only with three dots round every region ($6=3\cdot4-6$), so a fifth dot reaches only three of them. \emph{Limits:} counting proves impossibility, never possibility. The prism (network 4) has the same 6 dots and 9 links as $K_{3,3}$, and network 6 has 12 links, the most any untangled network on 6 dots can have, yet both untangle; only a drawing shows that. Kuratowski's theorem (not proved here) says a network untangles exactly when it contains no copy of $K_5$ or $K_{3,3}$, allowing extra dots along links.

\textbf{Fewest crossings.} $K_5$ and $K_{3,3}$ need exactly one crossing, and six dots all joined ($K_6$) need exactly three. A drawing gives each upper bound. For the lower bound, remove one link at each crossing: what is left is untangled, so with $k$ crossings $15-k\le 3\cdot6-6$, and $k\ge3$. This needs exactly two links at each crossing; if links could share one crossing place, all fifteen links of $K_6$ could pass through a single point.

\textbf{The link game is decided before it starts.} On $n\ge3$ dots every finished game has exactly $3n-6$ links however it is played, so the first player wins when $n$ is odd and the second when $n$ is even. \emph{Proof idea:} when no move is left, the drawing is in one piece (two pieces share a region, and a link through it joins them) and every region is a triangle. Otherwise some region has four or more dots round it, and two of them are unjoined: four dots round one region all joined to each other would let a new dot inside that region join all four, untangling $K_5$. A link through the region could then still be drawn. So $3F=2E$, and Euler's rule gives $E=3n-6$. \emph{Limits:} with straight links only, the count is $3n-3-h$, where $h$ dots lie on the outer boundary; every printed board has three outer dots, so straight and bent play agree there.

\textbf{The portal board is a torus.} On Week 41's board, with opposite sides glued, $K_{3,3}$, $K_5$, $K_6$ and even $K_7$ untangle (drawings at the printed positions are in the solutions). There $V-E+F=0$ when every region is one patch ($K_7$: $7-21+14=0$), so $E\le 3V$ and $K_8$ ($28>24$) cannot untangle even there. These are adult facts; the pages ask only for drawings.

\textbf{Experiments, conjectures, explanations.} Untangling attempts, the region table and game tallies are experiments; ``network 2 is stuck'', ``dots $-$ links $+$ regions $=2$'' and ``with 5 dots the first player wins'' are conjectures until explained. Many failed attempts are evidence, not proof; one good drawing proves that a network untangles. Every explanation above can be made with drawings and words.

\textbf{Bands.} K--1 untangles pipe-cleaner networks of four to six dots, finds that two houses and three wells can be joined without crossings while three and three get no better than one crossing, finds one crossing for five dots all joined, and plays the game with 4 and 5 dots; adults read. Grades 2--3 sort six networks into untangled and stuck, find Euler's rule, find the fewest crossings for networks 2 and 5 and for six dots, play the game, draw the stuck networks on the portal board, and end by explaining network 2. Grades 4--5 begin with the same sort, rule and game, prove networks 2 and 5 impossible, draw networks 5 and 2, then six and seven all-joined dots, on the portal board, prove six dots need three crossings, and explain the game for 10 and 11 dots.

\head{Preparation, launch and pacing}
\textbf{Status.} Prepared for organizer review; unpiloted. Nothing has been tried with real materials: paper dots with paper clips, pipe-cleaner ends twisted round the clips (four ends meet at every dot of network 2, network 6 and K--1 Problem 5), whether 30 cm pipe cleaners reach far enough when children spread the dots, whether a dot can be moved without its links coming off, and every time estimate below.

\textbf{Print.} US Letter, single-sided, at 100\% (Actual size), one packet per child: K--1 5 pages, Grades 2--3 7 pages, Grades 4--5 8 pages. Nothing is cut and there are no separate mats; the game boards (K--1 page 5, 2--3 page 5, 4--5 page 4) and portal boards (2--3 page 6, 4--5 pages 6 and 8) are packet pages. Print a few spares of those pages for second tries. Pencils with erasers.

\textbf{Kits.} One per pair, five in all (two pairs at K--1, two at 2--3, one at 4--5), plus a demonstration set of dots A--D and four pipe cleaners for the launch. A kit is six card dots about 5 cm across lettered A--F, each with a paper clip slid over its edge, and 15 pipe cleaners about 30 cm long; the printed networks need at most 6 dots and 12 links (network 6). Use letters only, no house or well pictures: C is a well in K--1 Problem 2 and a house in Problem 3, and children lay the dots out like the page. Round 5 cm stickers on card, or card squares, are quicker than hand-cut circles. The outline estimates about 20 minutes for five kits (not measured); with the demonstration set and printing, allow 30--40 minutes. Kits are reusable. If a crowded clip will not hold its ends, let the ends lie on the dot.

\textbf{Staffing.} One adult at each table. \emph{K--1, the parent volunteer, two pairs:} read each problem aloud, help twist ends onto clips, and copy a finished network into its box if a child wants; building and recording are ordinary support. The child chooses which dot to move; if the adult is choosing moves, go back to a smaller network. Before children try houses and wells, no adult says it is impossible; afterwards the adult may say that mathematicians have proved nobody can. \emph{2--3, the second mathematician, two pairs:} check region counts (expect the outside to be missed), show the game and read the portal rule at the table, and hold the explanation conversations without supplying the argument. \emph{4--5, the organizer, a pair and a third child:} the third child, in turn, checks each untangled drawing against its link list (every link present, no crossings) and referees the game, asking ``Is there a pair not joined that still can be?'' before a game is called. With only two adults, one can serve both older tables, whose children read for themselves.

\textbf{Launch (5 minutes, whole group, before the packets).} Let children handle the dots and pipe cleaners for two minutes, twisting one link onto two clips. Then, with children standing round one table, build this network, which is on no page: dots A, B, C, D; links A--B, B--C, C--A and C--D, with D beyond A--B so that C--D lies across it. Ask where two links cross, then move D; its link follows. Stop when no link lies on another: ``This one is untangled.'' Put D back and show a second legal way: leave D where it is and bend C--D round B. Show what is not allowed: a link lying across a dot, and unhooking a link. Draw the untangled network on the board, as children will draw theirs in the boxes. Do not untangle a page network, count regions, or mention houses and wells.\par
{\centering\fig{launch_in}{Input: 1 crossing}\hfill\fig{launch_step}{Move D; its link follows}\hfill\fig{launch_out}{Untangled}\hfill\fig{launch_bend}{Or bend C--D round B}\par}
\textbf{Two short demonstrations later, at each table.} When a table reaches the game, its adult and one child draw three links each on seven dots (whiteboard or scrap paper), and the adult points to one link that would cross and one that would join two dots already joined; stop long before the end and say nothing about how long a game lasts. When a table reaches the portal page, its adult traces the page's example with a finger (out at 1 on the right, back in at 1 on the left), then on an extra board draws one link from a dot out through the top and back in at the bottom.

\textbf{A first hour.} Running around and the launch take about 10 minutes, which leaves 35--40 minutes at the tables and about 3 minutes to close.\newline
K--1: Problem 1 (10--15 minutes; two networks are enough), Problems 2--3 (15 minutes), then the game, Problem 6 (10 minutes). Problems 4 and 5 suit quick pairs or a return visit. Children need not draw every network; draw Problems 2 and 3 and the best Problem 5.\newline
Grades 2--3: Problem 1 (20 minutes): one pair builds networks 1--3 and the other 4--6, then each tests the other's stuck network. Before Problem 2 the pairs show each other their untangled drawings and read out the dots, links and regions, so that every child's table has all four networks. Problem 2 (10--15 minutes), then the game, Problem 4, only if time allows. Return visit: the game if it was missed, then Problems 3, 5 and 6.\newline
Grades 4--5: Problem 1 (15--20 minutes; drawing is allowed), Problem 2 (10 minutes), the game, Problem 3 (10 minutes), and Problem 4 only if time allows. Return visit: Problems 4--9; a child stalled on a proof can go to the portal board, Problem 6.\newline
Close (3 minutes): each table shows one thing, such as K--1's best houses and wells, a stuck network from 2--3 and the 4--5 game tally. Finishing every page is not the goal.

\head{Grades K--1: untangle, then decide}
Accept any drawing with no crossings and the same links; below is one answer each, and for Problems 3 and 5 a drawing with the fewest crossings (circled). Ask ``How do you know for sure?'' only when a child is ready.
\begin{center}
\fig{key_k_bowtie}{P1, first}\hfill\fig{key_k_kite}{P1, second}\hfill\fig{key_k_k4}{P1, third}\hfill\fig{key_k_hw23}{P2}\hfill\fig{key_k_hw33}{P3: 1 crossing}\\[6pt]
\fig{key_k_star}{P4, first}\qquad\qquad\fig{key_k_fan}{P4, second}\qquad\qquad\fig{key_k_prism}{P4, third}\qquad\qquad\fig{key_k_k5}{P5: 1 crossing}
\end{center}

\sub{Problem 1, page 1: three four-dot networks}
All three untangle. Each picture has one crossing, and moving one dot fixes it. The ring A--B--C--D becomes a square; in the second, A and D go on opposite sides of link B--C; in the third, where every pair is joined, any one dot goes inside the triangle of the other three. Hint: ``Which two links cross? What if one of their dots moves?'' Extension for the third: ``Can you add one more link?'' (No; every pair is already joined.)

\sub{Problems 2 and 3, page 2: houses and wells}
Two houses and three wells: yes, with one house above the row of wells and the other below it. Hint: ``What if one house moves to the other side of the wells?'' Three houses and three wells: no; one crossing is the fewest (key, P3). Expect many children to get down to one crossing and stop. ``We couldn't, and one crossing is the fewest we found'' is the expected K--1 answer; it is evidence, not a proof, and the 4--5 pages prove it. Hint: ``How few crossings can you get?'' Extension: ``Take away any one pipe cleaner. Can you untangle it now?'' (Yes, whichever one goes.)

\sub{Problem 4, page 3: five and six dots}
All three untangle; the pictures have 5, 3 and 4 crossings. The first is one ring, A--B--C--D--E. In the second, E is joined to every other dot and A--B--C--D is a path, so E goes above a row A, B, C, D. The third is two triangles, A, B, C and D, E, F, joined by A--D, B--E and C--F; one triangle goes inside the other. Hints: ``Start at A and follow the links. Which dot comes next?'' ``Can you find three dots all joined to each other?''

\sub{Problem 5, page 4: five dots, all joined}
One crossing is the fewest (key, P5); expect children to reach two or three first. Ten links is a lot of twisting: an adult may attach them while the child decides where the dots go. An argument an adult can show with the pipe cleaners: untangle A, B, C, D first, as in Problem 1's third network. Every region then has three dots round it, so wherever E goes it reaches only three of the four without a crossing. Hint: ``Can you untangle four of the dots first? Where can the fifth go?''

\sub{Problem 6, page 5: the link game}
With 4 dots the second player wins every game, and with 5 dots the first player does, however anyone plays: a finished game has 6 links on 4 dots and 9 on 5. On 4 dots it joins every pair; on 5 dots exactly one pair stays unjoined (compare Problem 5). The risk is stopping early, when a link between the inner dots or along the outer triangle still fits. Before a game ends, ask: ``Is there a pair of dots that is not joined and still can be?'' If everyone draws straight links, the game is over exactly when no straight link fits; once someone draws a curved link, check curved routes too, including round the outside. Either way the total is the same. Let each child play first in some games and second in others, so the tally shows that the winning seat, not the player, repeats. Extension: ``Count the links at the end of each game.''

\head{Grades 2--3: sort, count, and the portal board}
\begin{center}
\fig{key_n1}{Network 1: 5, 7, 4}\hfill\fig{key_n3}{Network 3: 5, 9, 6}\hfill\fig{key_n4}{Network 4: 6, 9, 5}\hfill\fig{key_n6}{Network 6: 6, 12, 8}\\[4pt]
{\small\textbf{Figure 1.} The networks that untangle, with dots, links and regions.}\\[10pt]
\fig{few_K5}{Network 2: 1 crossing}\qquad\fig{few_K33}{Network 5: 1 crossing}\qquad\fig{few_K6}{6 dots: 3 crossings}\\[4pt]
{\small\textbf{Figure 2.} Fewest crossings, circled. Network 2's drawing is network 3's plus D--E.}
\end{center}

\sub{Problem 1, pages 1--2: untangled or stuck}
Networks 1, 3, 4 and 6 untangle (Figure 1); networks 2 and 5 stay stuck. The pictures have 2, 3, 3, 4, 5 and 4 crossings. The pairs are deliberate: network 3 is network 2 without D--E; networks 4 and 5 both have 6 dots and 9 links; network 6 has the most links and still untangles. Network 5 is the houses and wells. At this point ``stuck'' is a conjecture. Problem 6 asks why for network 2; network 5 stays a conjecture in this packet, though a ready child can try the overview's ring argument. Hints: for network 6, ``Which dot is A not joined to?'' (F; B--D and C--E are also missing) and ``What if A, B and C made a big triangle?'' For network 3, ``Which pair of dots is not joined?''

\sub{Problem 2, page 3: dots, links and regions}
Dots $-$ links $+$ regions $=2$ for every untangled network in one piece; children may write ``regions $=$ links $-$ dots $+2$'' or ``dots $+$ regions $=$ links $+2$''. The rows are in Figure 1, and the example has 4, 5, 3. Likely slips: forgetting the outside, counting a picture that still has crossings, and losing track in a crowded drawing; shading each region as it is counted helps. Hint: ``What happens with just one link? With two?'' Extensions: test a network in two pieces (it gives 3), and count the regions of a tangled picture (each crossing adds one). For a child who asks why: add the links one at a time and watch the three numbers.

\sub{Problem 3, page 4: as few crossings as you can}
Network 2: 1. Network 5: 1. Six dots all joined: 3 (Figure 2). Untangle network 3 and add D--E, which crosses one link; untangle houses and wells with one link removed and put it back across one link. For six dots, try one triangle inside another. That 1 is the best for network 2 is Problem 6's question. That 1 is the best for network 5 and 3 for six dots stay conjectures here (4--5 Problems 5 and 7 prove them). Hint: ``Which link could you take away to untangle it?''

\sub{Problem 4, page 5: the link game}
4 dots: second player; 5 dots: first player; 6 dots: second player, in every game. Finished games have 6, 9 and 12 links. Check endings as in K--1 Problem 6: ``Is there a pair that is not joined and still can be?'' Extensions: ``Count the links at the end.'' ``Predict 7 dots, then test.'' (15 links; the first player wins.)

\begin{center}
\fig{portal_k33}{Network 5}\qquad\qquad\fig{portal_k5}{Network 2}\\[4pt]
{\small\textbf{Figure 3.} One portal-board drawing at the printed positions. As in the page's example, a link that leaves at a number comes back in at the same number on the opposite side.}
\end{center}

\sub{Problem 5, page 6: the portal board}
Both can be drawn (Figure 3). Network 5: A joins all three wells directly, as do B--F and C--F; B--D, B--E and C--E leave through the top and come back in at the bottom; C--D leaves through the right side and comes back on the left. Network 2: join each pair by its shortest route, counting routes through the sides; five links use a portal. Many other drawings work. Hints: ``Which link is blocked inside the board? Could it leave through a side?'' ``Where does it come back in?'' Extension: count regions on the portal board, following each region across the sides. Every crossing-free drawing gives 3 for network 5 and 5 for network 2, so dots $-$ links $+$ regions $=0$ here.

\sub{Problem 6, page 7: why network 2 is stuck}
Three explanations children can reach. (1) A, B, C and D are all joined. Untangled, every region has just three of these dots round it (one dot sits inside the triangle of the other three; in general $6=3\cdot4-6$, so every region is a triangle), and E, in whichever region it sits, reaches only those three. (2) Draw the ring A--B--C--D--E first; each of the other five links goes inside or outside it, two fit inside and two outside, and the fifth must cross. (3) Untangled, Problem 2's rule gives 7 regions, each needing at least 3 links round it: 21 link-sides, but 10 links have 20. ``We tried lots of ways'' is evidence; ask ``Could someone else find a way? How do you know?'' Accept a partial argument with a drawing.

\head{Grades 4--5: proofs, the game and the portal board}
\sub{Problems 1 and 2, pages 1--3}
As Grades 2--3 Problems 1 and 2, without the split: the pair does all six networks and may draw instead of building.

\sub{Problem 3, page 4: the link game}
As Grades 2--3 Problem 4: second, first, second player. Collect the tally now; Problem 8 asks why.

\sub{Problems 4 and 5, page 5: why networks 2 and 5 are stuck}
Network 2: if it were untangled, Problem 2's rule would give $10-5+2=7$ regions. With no double links, each region has at least 3 links round it, so at least 21 link-sides, but 10 links have only 20. The ring and four-dot arguments of Grades 2--3 Problem 6 are also complete proofs. Follow-up: ``Why at least three links round a region?'' (Two would mean a double link.)
Network 5: every link joins a house to a well, so a loop alternates houses and wells and needs at least 4 links; no three dots are all joined. Untangled: $9-6+2=5$ regions, at least 4 link-sides each, so 20, but 9 links have 18. Ring version: around A--D--B--E--C--F the links A--E, B--F and C--D join opposite corners, any two on the same side cross, and three links cannot share two sides. Hint: ``Can you find a triangle in network 5? What does that say about its regions?''

\sub{Problem 6, page 6: the portal board}
As Grades 2--3 Problem 5.

\sub{Problems 7 and 8, page 7: six dots, then 10 and 11 dots}
Problem 7: 3 crossings (Figure 2). Nobody can do better: with 2 crossings, removing one link at each would leave at least 13 links on 6 dots untangled, but an untangled network on 6 dots has at most 12 (from regions: $3F\le 2E$ and $F=E-4$ give $E\le 12$; or because it can be continued as a game, which ends at 12). Only two links meet at a crossing, which is why one removal per crossing is enough. Follow-up: ``Could one link be in two crossings?'' (Removing it then clears both, which only helps.) Hint: ``What if you take away one link at each crossing?''

Problem 8: the second player wins with 10 dots (24 links) and the first with 11 (27). When nobody can move, the network is in one piece and every region is a triangle, or another link would fit; so $3\times\text{regions}=2\times\text{links}$, and with $\text{dots}-\text{links}+\text{regions}=2$ this gives $\text{links}=3\times\text{dots}-6$. Expect children to see 6, 9, 12 and ``three more links per dot'' first. That is a conjecture: putting a new dot inside a triangle of a finished game and joining it to the corners makes one finished game with three more links, not every game. Hint: ``When nobody can move, what do the regions look like?'' Follow-up: ``Why can't a finished game have a region with four dots round it?'' (If P, Q, R, S are in order round it, P--R and Q--S cannot both run outside the region without crossing, so one of them can still be drawn inside.)

\sub{Problem 9, page 8: six and seven dots on the portal board}
Both work at the printed positions (Figure 4): join each pair by its shortest route through the repeated board, and no two links cross. That is a property of these positions, not a general rule; for network 5 the shortest routes tie. Seven dots: 21 links, 10 through a portal, and the board fills with 14 triangles. Six dots: the same drawing without G, 15 links, 8 through a portal, 9 regions. Hint, only when needed: ``If the board repeated in every direction, as on Week 41's map, which copy of each dot is nearest?'' Extension for adults to discuss: eight dots all joined cannot fit even here, since $28>3\times8$.\par\needspace{8cm}
\begin{center}
\fig{portal_k6}{6 dots}\qquad\qquad\fig{portal_k7}{7 dots}\\[4pt]
{\small\textbf{Figure 4.} Six and seven dots at the printed positions, each pair joined by its shortest route; numbers as in Figure 3.}
\end{center}

\head{Sources, connections and checks}
\textbf{Mathematics.} Douglas B. West, \emph{Introduction to Graph Theory}, 2nd ed. (Prentice Hall, 2001), \S6.1: Euler's formula and $E\le3V-6$, $E\le2V-4$; section number from the outline, not consulted for this guide. K. Kuratowski, ``Sur le probl\`eme des courbes gauches en topologie'', \emph{Fundamenta Mathematicae} 15 (1930): the obstruction theorem, stated only; not consulted. The proofs here, the game analysis and the portal drawings are written for this guide and checked by computation.

\textbf{Puzzles.} Henry E. Dudeney, \emph{Amusements in Mathematics} (1917), ``Water, Gas, and Electricity'': problem 251, p.~73, solution pp.~200--201, as cited by the Wikipedia page ``Three utilities problem'' (consulted October 25, 2026); the book itself was not checked. That page reports that Dudeney's solution passes a pipe through a house, the loophole closed here by ``never lies across a dot''. ``Planarity'', Wikipedia (consulted October 25, 2026): John Tantalo's 2005 game, from a concept by Mary Radcliffe, in which players move dots one at a time until no links cross; the untangling format comes from it. Simon Tatham, ``Untangle'' (consulted October 25, 2026): the same puzzle, credited to Planarity.

\textbf{Library.} Week 56 proves $V-E+F=2$ on solids; here it becomes a tool for proving a drawing impossible. Week 41 built the portal board, whose seam marks these pages reuse. Week 40 studies crossings in knot drawings, and Week 10 counts the links at each island.

\textbf{Checks and what to record.} \texttt{check\_guide.py} recomputes every answer above from a rebuild of the final pages, including 300 random games on each printed game board, the bent launch link and the portal drawings of Figures 3 and 4 with their numbered exits; none of this is a physical rehearsal. Record the real kit time, whether clips hold four pipe-cleaner ends, whether games stop early, which networks children called stuck and why, the region rule in children's words, and which problems children want to continue.

\end{document}
